\section{Presupuesto}

Precios de 2026 en pesos colombianos. Los materiales los financia el estudiante; el personal se valora pero no representa desembolso. La Universidad aporta en especie taller, laboratorios e instrumentación.
\vspace{-4pt}

\begin{table}[H]
\centering
\singlespacing
\scriptsize
\caption{Presupuesto estimado del proyecto (COP, 2026).}
\label{tab:presupuesto}
\renewcommand{\arraystretch}{0.95}
\begin{tabular}{p{8.6cm} c r}
\toprule
\textbf{Rubro} & \textbf{Fuente} & \textbf{Valor (COP)}\\
\midrule
Estudiante: \num{384} h (\num{16} semanas $\times$ \num{24} h) a \cop{25000}/h & Especie & \cop{9600000}\\
Asesor: \num{32} h a \cop{90000}/h & EAFIT & \cop{2880000}\\
\midrule
Subsistema térmico: resistencia, canal reflector, aislamiento y termopares & Estudiante & \cop{280000}\\
Control, interfaz y accionamiento: Arduino, SSR, HMI táctil y motor NEMA 23 & Estudiante & \cop{880000}\\
Estructura y mecanizado: perfilería, lámina, tornillería y transmisión & Estudiante & \cop{950000}\\
Fuente, gabinete, cableado, protecciones y paro de emergencia & Estudiante & \cop{300000}\\
Acrílico de ensayo, instrumentación de validación e imprevistos & Estudiante & \cop{736000}\\
\midrule
\textbf{Subtotal personal} \, / \, \textbf{materiales e insumos} & & \textbf{\cop{12480000}} \, / \, \textbf{\cop{3146000}}\\
\textbf{Total} & & \textbf{\cop{15626000}}\\
\bottomrule
\end{tabular}
\end{table}
